\chapter{From Cooper-pair box to transmon}
\label{ch:transmon}

The Josephson Hamiltonian
$\op H=4E_C(\op n-n_g)^2-E_J\cos\op\varphi$ of
\cref{stmt:josephson-hamiltonian} contains a whole family of qubits;
which one it describes is decided entirely by the ratio $E_J/E_C$. At
small $E_J/E_C$ the charging term dominates: charge is nearly a good
quantum number, the eigenstates are Cooper-pair-number states weakly
mixed by tunnelling, and the device is the \emph{Cooper-pair box} ---
the first superconducting qubit, exquisitely anharmonic but ruined by
charge noise. At large $E_J/E_C$ the Josephson term dominates: the
phase is pinned near the bottom of a cosine well, the circuit is a
weakly anharmonic oscillator, and the device is the \emph{transmon},
which trades most of the box's anharmonicity for an exponential
immunity to charge noise. This chapter diagonalises the Hamiltonian in
both limits, extracts the transmon's transition frequency
$\hbar\omega_{01}=\sqrt{8E_JE_C}-E_C$ and its negative anharmonicity
$\alpha\approx-E_C$, explains why charge dispersion vanishes
exponentially in $\sqrt{E_J/E_C}$, and closes with a brief tour of the
other corners of the design landscape --- the flux qubit and
fluxonium. The output is the concrete qubit whose abstract
two-level version animated Parts~I--III, ready to be coupled to a
resonator in \cref{ch:circuit-qed}.

% --------------------------------------------------------------
\section{Anharmonicity makes a qubit}
\label{sec:anharmonicity}

\begin{phenomenon}
A qubit is a \emph{two}-level system, but a Josephson circuit has
infinitely many levels. What lets us use only the lowest two is that
the level spacing is unequal: the energy to go from $\ket0$ to $\ket1$
differs from the energy to go from $\ket1$ to $\ket2$. A drive resonant
with the first transition is then off-resonant with the second, so it
shuffles population only within $\{\ket0,\ket1\}$ and leaves the higher
levels alone. The larger this mismatch, the faster one may drive
without leaking out of the computational subspace. Anharmonicity is the
resource that turns an oscillator into a qubit, and its size sets the
speed limit on gates.
\end{phenomenon}

\begin{statement}[\textbf{Anharmonicity and the gate-speed limit}]
\label{stmt:anharmonicity}
Let $\hbar\omega_{ij}=E_j-E_i$ be the transition energies of a circuit.
The \emph{anharmonicity} is
\begin{equation}
\label{eq:anharmonicity}
\alpha \;\equiv\; \hbar\omega_{12}-\hbar\omega_{01}
= (E_2-E_1)-(E_1-E_0),
\end{equation}
the deviation of the spectrum from the uniform ladder of the linear
resonator (\cref{stmt:lc-oscillator}), for which $\alpha=0$. To address
$\ket0\!\leftrightarrow\!\ket1$ without exciting
$\ket1\!\leftrightarrow\!\ket2$, a control pulse of duration $\tau$ ---
hence spectral width $\sim2\pi/\tau$ --- must satisfy
$|\alpha|/\hbar\gg2\pi/\tau$, so the fastest leakage-free gate scales
as $\tau\gtrsim2\pi\hbar/|\alpha|$. A finite, sufficiently large
$|\alpha|$ is therefore the single requirement that the bare cosine of
\cref{ch:josephson} must meet.
\end{statement}

\paragraph{Transition.}
Everything now hinges on how the spectrum of
\cref{eq:josephson-hamiltonian} deviates from uniform. We compute it in
the two opposite limits of $E_J/E_C$, beginning with the
charge-dominated Cooper-pair box.

% --------------------------------------------------------------
\section{The Cooper-pair box}
\label{sec:cpb}

\begin{phenomenon}
Isolate a small superconducting island connected to a reservoir by a
single junction and a gate capacitor. So few Cooper pairs sit on the
island that the energy to add one more, the charging energy $E_C$,
dominates. The natural states are then definite charge states
$\ket n$ --- $n$ excess Cooper pairs --- and the Josephson coupling
acts only as a weak amplitude to move one pair on or off, mixing
neighbouring charge states. A gate voltage tunes the electrostatic
energy of each $\ket n$, and at the special bias where two charge
states cost the same energy the junction lifts their degeneracy: there
sits a clean, highly anharmonic two-level system.
\end{phenomenon}

\begin{statement}[\textbf{Cooper-pair box in the charge basis}~\cite{nakamura1999}]
\label{stmt:cpb}
In the charge basis $\{\ket n\}$ (eigenstates of $\op n$), the
Josephson term is a nearest-neighbour hopping, since
$\cos\op\varphi=\tfrac12\sum_n(\ket{n}\!\bra{n+1}+\ket{n+1}\!\bra{n})$,
so the Hamiltonian \cref{eq:josephson-hamiltonian} reads
\begin{equation}
\label{eq:cpb-hamiltonian}
\op H = \sum_n\Bigl[4E_C(n-n_g)^2\,\ket n\!\bra n
  - \frac{E_J}{2}\bigl(\ket{n}\!\bra{n+1}+\ket{n+1}\!\bra{n}\bigr)\Bigr].
\end{equation}
In the charge regime $E_J\ll E_C$, near the degeneracy point
$n_g=\tfrac12$ only $\ket0$ and $\ket1$ matter; projected onto them,
\begin{equation}
\label{eq:cpb-twolevel}
\op H_{\rm 2L} = 2E_C(1-2n_g)\,\hat\sigma_z - \frac{E_J}{2}\hat\sigma_x,
\end{equation}
a spin in a field. At the \emph{sweet spot} $n_g=\tfrac12$ the
longitudinal field vanishes and the gap is $E_J$, set purely by the
Josephson energy; the two eigenstates are the symmetric/antisymmetric
combinations $(\ket0\pm\ket1)/\sqrt2$.
\end{statement}

\paragraph{Charge dispersion and its curse.}
Away from the sweet spot the levels disperse with the offset charge
$n_g$ --- the \emph{charge dispersion}, a set of $n_g$-periodic energy
bands of width $\sim E_J$. Because real devices suffer slow,
uncontrolled drifts of $n_g$ (charges hopping in the substrate and
oxide), the qubit frequency wanders, and the qubit dephases through the
longitudinal noise mechanism of \cref{stmt:noise-rates}. Operating
exactly at the sweet spot $n_g=\tfrac12$ makes the frequency
first-order insensitive to $n_g$ --- the same protection idea as the
SQUID flux sweet spot (\cref{stmt:squid}) --- and Nakamura's first
charge qubit~\cite{nakamura1999} demonstrated coherent control there.
But the protection is only first order: charge noise still dephases the
box on nanosecond-to-microsecond timescales, and one cannot stay
pinned at the sweet spot during arbitrary gates. This was the wall the
field hit, and the transmon is how it was climbed.

\paragraph{Transition.}
The cure is counter-intuitive: rather than fight charge noise at fixed
$E_J/E_C$, simply increase $E_J/E_C$ until the entire charge dispersion
--- the whole $n_g$-dependence --- shrinks below the noise. We now show
it shrinks \emph{exponentially}.

% --------------------------------------------------------------
\section{The transmon limit}
\label{sec:transmon-limit}

\begin{phenomenon}
Shunt the junction with a large capacitor, lowering $E_C$ until
$E_J/E_C$ reaches several tens. The phase $\varphi$ is now heavy and
sits deep in a single cosine well, oscillating with small amplitude
about the minimum. The circuit becomes a slightly anharmonic
oscillator: its lowest levels are nearly those of a harmonic oscillator
at the plasma frequency, with a small negative correction from the
quartic curvature of the cosine. Crucially, the wavefunctions are now
broad in phase and therefore narrow in charge --- so spread over many
charge states that they no longer resolve the offset $n_g$, and charge
noise becomes invisible. The modest price is that the anharmonicity,
once the full $E_J$ of the box, drops to about $E_C$.
\end{phenomenon}

\begin{statement}[\textbf{Transmon spectrum}~\cite{koch2007}]
\label{stmt:transmon}
For $E_J/E_C\gg1$, expanding $-E_J\cos\op\varphi$ about $\varphi=0$ and
treating the quartic term in first-order perturbation theory gives the
levels
\begin{equation}
\label{eq:transmon-levels}
E_m \approx -E_J + \sqrt{8E_JE_C}\,\Bigl(m+\tfrac12\Bigr)
  - \frac{E_C}{12}\bigl(6m^2+6m+3\bigr),
\end{equation}
so that the transition frequency, anharmonicity, and their ratio are
\begin{equation}
\label{eq:transmon-freq}
\hbar\omega_{01}=\sqrt{8E_JE_C}-E_C,
\qquad
\alpha=\hbar\omega_{12}-\hbar\omega_{01}=-E_C,
\qquad
\frac{\alpha}{\hbar\omega_{01}}\approx-\sqrt{\frac{E_C}{8E_J}}.
\end{equation}
The charge dispersion --- the peak-to-peak variation of $E_m$ with the
offset $n_g$ --- is suppressed exponentially,
\begin{equation}
\label{eq:charge-dispersion}
\epsilon_m \;\sim\; E_C\,e^{-\sqrt{8E_J/E_C}},
\end{equation}
(times an $m$-dependent algebraic prefactor that grows with the level),
so charge noise is rendered negligible for the lowest levels once
$E_J/E_C\gtrsim50$.
\end{statement}

\paragraph{Derivation (Duffing oscillator).}
Expand $-E_J\cos\op\varphi=-E_J+\tfrac12 E_J\op\varphi^2
-\tfrac1{24}E_J\op\varphi^4+\cdots$. The quadratic piece together with
$4E_C\op n^2$ is a harmonic oscillator: with $\comm{\op\varphi}{\op n}=i$,
identifying $\op\varphi$ as ``position'' and $\op n$ as ``momentum''
with effective mass $1/8E_C$ and spring constant $E_J$ gives plasma
frequency $\hbar\omega_p=\sqrt{8E_JE_C}$ and zero-point phase spread
$\varphi_{\rm zpf}=(2E_C/E_J)^{1/4}$, so
$\op\varphi=\varphi_{\rm zpf}(\hat a+\hat a^\dagger)$. The quartic term
becomes $-\tfrac{1}{24}E_J\varphi_{\rm zpf}^4(\hat a+\hat a^\dagger)^4$;
its first-order (normal-ordered, RWA) contribution is a \emph{Kerr}
nonlinearity
\begin{equation}
\label{eq:transmon-kerr}
\op H_{\rm tr}\approx \hbar\tilde\omega\,\hat a^\dagger\hat a
  - \frac{E_C}{2}\,\hat a^\dagger\hat a^\dagger\hat a\hat a,
\end{equation}
with $\hbar\tilde\omega=\sqrt{8E_JE_C}-E_C$. The anharmonic shift of
level $m$ is $-\tfrac{E_C}{2}m(m-1)$, reproducing
\cref{eq:transmon-levels} and giving the spacing
$\hbar\omega_{m,m+1}=\hbar\tilde\omega-mE_C$, hence $\alpha=-E_C$. The
charge dispersion \cref{eq:charge-dispersion} follows from a WKB
estimate of the tunnelling amplitude between adjacent cosine wells: it
is the bandwidth of a particle of mass $1/8E_C$ in a periodic
potential of depth $E_J$, exponentially small in
$\sqrt{E_J/E_C}$.

\begin{remark}
The transmon is thus a \emph{Kerr oscillator}: a harmonic oscillator
with a weak photon-photon interaction $-\tfrac{E_C}{2}\hat a^\dagger\hat a^\dagger\hat a\hat a$.
The negative sign means each added excitation lowers the next
transition frequency by $E_C$ --- the same negative anharmonicity that
entered the transmon correction to the dispersive shift in
\cref{stmt:dispersive-H} (the factor $\alpha/(\Delta+\alpha)$). The
Kerr form is also why a transmon, driven hard, behaves like the
anharmonic oscillators that generate the parametric and mixing
processes of \cref{ch:circuit-qed}.
\end{remark}

\paragraph{The exponential trade-off.}
\Cref{eq:transmon-freq,eq:charge-dispersion} encode the central
compromise. Charge dispersion falls as $e^{-\sqrt{8E_J/E_C}}$ ---
\emph{exponentially} fast --- while the anharmonicity falls only as
$\alpha/\omega_{01}\sim-\sqrt{E_C/8E_J}$ --- a gentle \emph{power law}.
Increasing $E_J/E_C$ therefore buys exponential charge-noise immunity
at merely polynomial cost in anharmonicity. The standard operating
point $E_J/E_C\approx50$--$100$ leaves $\epsilon_m$ utterly negligible
(qubit frequency stable to $\ll1\,$Hz from charge noise) while keeping
$|\alpha|/h\approx E_C/h\approx200$--$300\,$MHz --- enough for
$\sim20$--$50\,$ns gates. That asymmetry is the whole reason the
transmon won.

\paragraph{Transition.}
The transmon is the dominant qubit, but not the only one. Two further
designs occupy different corners of the $E_J/E_C$ (and inductive-energy)
landscape, trading the transmon's balance for larger anharmonicity.

% --------------------------------------------------------------
\section{Flux qubit and fluxonium}
\label{sec:flux-fluxonium}

\begin{statement}[\textbf{Two other corners of the design space}]
\label{stmt:flux-fluxonium}
\textbf{Flux qubit.} A superconducting loop interrupted by (typically
three) junctions and biased near half a flux quantum has two low-lying
states corresponding to persistent currents circulating clockwise and
anticlockwise. Near $\Phi_{\rm ext}=\Phi_0/2$ these states are
degenerate and the junctions tunnel-couple them, giving a two-level
system with \emph{large} anharmonicity but strong sensitivity to flux
noise away from the half-flux sweet spot.

\textbf{Fluxonium.} A single small junction shunted by a very large
inductance --- a ``superinductance'' built from an array of larger
junctions --- and a small capacitance~\cite{manucharyan2009}. With
inductive energy $E_L\ll E_J$, the cosine potential is corrugated by a
shallow parabola, producing a dense set of nearly flux-independent
levels with very large anharmonicity and, in the heavy-fluxonium
regime, exceptionally long coherence times. Like the transmon, it is
free of charge offsets.
\end{statement}

\paragraph{Why the transmon still dominates.}
Each alternative buys anharmonicity at a cost: the flux qubit is
flux-noise-limited; the fluxonium needs a superinductance and operates
at low frequency, complicating control and readout. The transmon's
moderate $|\alpha|\approx E_C$ is ``good enough'' for fast gates, while
its exponential charge-noise immunity, simple single-junction (or
SQUID) construction, and convenient $4$--$8\,$GHz frequencies make it
the default. The rest of Part~IV therefore takes the transmon as the
qubit, and reads its coupling, control, and calibration off the cQED
machinery of \cref{ch:circuit-qed,ch:experimental-basics}.

% --------------------------------------------------------------
This chapter diagonalised the Josephson Hamiltonian and met the
superconducting qubit in two guises. The Cooper-pair box, in the
charge regime $E_J\ll E_C$, is a clean two-level system at its charge
sweet spot but is dephased by ubiquitous charge noise. The transmon,
in the regime $E_J/E_C\gg1$, suppresses charge dispersion
exponentially in $\sqrt{E_J/E_C}$ at only a power-law cost in
anharmonicity, leaving a weakly anharmonic Kerr oscillator with
$\hbar\omega_{01}=\sqrt{8E_JE_C}-E_C$ and $\alpha\approx-E_C$. Flux
qubit and fluxonium occupy other corners of the landscape, trading the
transmon's balance for larger anharmonicity.

With a concrete qubit in hand we can finally give the abstract objects
of Parts~I--III their hardware. The next chapter places a transmon in a
microwave resonator, derives the Jaynes--Cummings coupling $g$ from
circuit capacitances, builds single- and two-qubit gates from drives
and tunable couplings, and shows how the same Josephson nonlinearity
generates the parametric processes behind quantum-limited
amplification.

\clearpage
\begin{chaptersummary}

\paragraph{Anharmonicity.}
$\alpha=\hbar\omega_{12}-\hbar\omega_{01}$ measures departure from the
uniform ladder; leakage-free gates need $|\alpha|/\hbar\gg2\pi/\tau$,
so $\tau\gtrsim2\pi\hbar/|\alpha|$.

\paragraph{Cooper-pair box ($E_J\ll E_C$).}
Charge-basis Hamiltonian
$\op H=\sum_n[4E_C(n-n_g)^2\ket n\!\bra n-\tfrac{E_J}{2}(\ket n\!\bra{n+1}+\hc)]$;
near $n_g=\tfrac12$ a spin
$\op H_{\rm 2L}=2E_C(1-2n_g)\hat\sigma_z-\tfrac{E_J}{2}\hat\sigma_x$,
gap $E_J$ at the sweet spot. Large anharmonicity, but charge-noise
limited.

\paragraph{Transmon ($E_J/E_C\gg1$).}
$\hbar\omega_{01}=\sqrt{8E_JE_C}-E_C$, $\alpha=-E_C$,
$\alpha/\hbar\omega_{01}\approx-\sqrt{E_C/8E_J}$. Charge dispersion
$\epsilon_m\sim E_Ce^{-\sqrt{8E_J/E_C}}$ is exponentially suppressed; a
Kerr oscillator
$\op H\approx\hbar\tilde\omega\hat a^\dagger\hat a-\tfrac{E_C}{2}\hat a^\dagger\hat a^\dagger\hat a\hat a$.
Operate at $E_J/E_C\approx50$--$100$.

\paragraph{Other designs.}
Flux qubit (persistent-current states near $\Phi_0/2$, large $\alpha$,
flux-noise-limited) and fluxonium (superinductance shunt, very large
$\alpha$, long coherence) trade the transmon's balance for
anharmonicity.

\end{chaptersummary}

% --------------------------------------------------------------
\section*{Exercises}
\addcontentsline{toc}{section}{Exercises}

\begin{exercise}[The Josephson term as charge hopping]
Show that $\cos\op\varphi$ moves one Cooper pair on or off the island,
i.e.\ that
$e^{\pm i\op\varphi}\ket n=\ket{n\pm1}$, and hence justify the
nearest-neighbour form in \cref{eq:cpb-hamiltonian}.
\begin{solution}
Since $\comm{\op\varphi}{\op n}=i$, the operator $e^{i\op\varphi}$ is
the translation operator in the charge variable conjugate to
$\varphi$: $e^{i\op\varphi}\ket n=\ket{n+1}$ (and $e^{-i\op\varphi}\ket n=\ket{n-1}$),
exactly as $e^{ipa/\hbar}$ shifts position. Then
$\cos\op\varphi=\tfrac12(e^{i\op\varphi}+e^{-i\op\varphi})$ maps
$\ket n\to\tfrac12(\ket{n+1}+\ket{n-1})$, i.e.\ in matrix form
$\cos\op\varphi=\tfrac12\sum_n(\ket n\!\bra{n+1}+\ket{n+1}\!\bra n)$,
the nearest-neighbour hopping of \cref{eq:cpb-hamiltonian}.
\end{solution}
\end{exercise}

\begin{exercise}[The box at its sweet spot]
Diagonalise \cref{eq:cpb-twolevel} at $n_g=\tfrac12$ and find the gap
and eigenstates. What is the gap a few percent away from the sweet
spot?
\begin{solution}
At $n_g=\tfrac12$ the $\hat\sigma_z$ term vanishes and
$\op H_{\rm 2L}=-\tfrac{E_J}{2}\hat\sigma_x$, with eigenvalues
$\mp E_J/2$ and eigenstates $(\ket0\pm\ket1)/\sqrt2$; the gap is $E_J$.
Away from the sweet spot, $n_g=\tfrac12+\delta$, the field is
$\sqrt{(4E_C\cdot2\delta)^2+E_J^2}=\sqrt{(8E_C\delta)^2+E_J^2}$, so to
leading order the gap grows quadratically in $\delta$ --- the
first-order insensitivity that defines the sweet spot.
\end{solution}
\end{exercise}

\begin{exercise}[Transmon frequency and the $\sqrt{8E_JE_C}$ scale]
A transmon has $E_J/h=15\,$GHz and $E_C/h=0.25\,$GHz. Compute
$\omega_{01}/2\pi$ and the anharmonicity $\alpha/h$.
\begin{solution}
$\sqrt{8E_JE_C}=\sqrt{8\cdot15\cdot0.25}\,$GHz$\cdot h
=\sqrt{30}\,$GHz$\approx5.48\,$GHz (working in frequency units, since
the $h$'s cancel). Then
$\omega_{01}/2\pi=\sqrt{8E_JE_C}/h-E_C/h\approx5.48-0.25=5.23\,$GHz, and
$\alpha/h=-E_C/h=-0.25\,$GHz$=-250\,$MHz. Relative anharmonicity
$\alpha/\hbar\omega_{01}\approx-250/5230\approx-4.8\%$.
\end{solution}
\end{exercise}

\begin{exercise}[Choosing $E_J/E_C$]
Using \cref{eq:charge-dispersion}, estimate by what factor the charge
dispersion changes when $E_J/E_C$ goes from $20$ to $80$. Comment on
the corresponding change in anharmonicity.
\begin{solution}
$\epsilon\propto e^{-\sqrt{8E_J/E_C}}$. At $E_J/E_C=20$,
$\sqrt{8\cdot20}=\sqrt{160}\approx12.6$; at $80$,
$\sqrt{640}\approx25.3$. The ratio of dispersions is
$e^{-(25.3-12.6)}=e^{-12.7}\approx3\times10^{-6}$ --- six orders of
magnitude smaller. The relative anharmonicity
$\sqrt{E_C/8E_J}$ changes only by $\sqrt{80/20}=2$ (halved). Charge
noise is crushed exponentially; anharmonicity is sacrificed only by a
factor of two --- the bargain that defines the transmon.
\end{solution}
\end{exercise}

\begin{exercise}[Zero-point phase spread]
Compute $\varphi_{\rm zpf}=(2E_C/E_J)^{1/4}$ for $E_J/E_C=60$ and
confirm it is small ($\ll1$), justifying the expansion of the cosine.
\begin{solution}
$\varphi_{\rm zpf}=(2/60)^{1/4}=(0.0333)^{1/4}$. Since
$0.0333^{1/2}=0.183$ and $0.183^{1/2}=0.427$, $\varphi_{\rm zpf}\approx0.43\,$rad.
This is well below $1$, so the phase explores only the bottom of the
cosine well and the quartic-truncated (Duffing) treatment is justified;
the next, sextic term is suppressed by a further
$\varphi_{\rm zpf}^2\approx0.18$.
\end{solution}
\end{exercise}

\begin{exercise}[Anharmonicity sets the gate speed]
With $|\alpha|/h=250\,$MHz, estimate the shortest gate time before
leakage to $\ket2$ becomes significant. Compare to typical $T_1\sim
50\,\mu$s.
\begin{solution}
Leakage-free operation needs the pulse bandwidth
$\sim1/\tau$ below $|\alpha|/h$, so
$\tau\gtrsim1/(|\alpha|/h)=1/(250\times10^6\,\text{Hz})=4\,$ns as an
absolute floor; in practice a safety factor of $\sim5$--$10$ (and pulse
shaping, e.g.\ DRAG) gives $\sim20$--$40\,$ns gates. Against
$T_1\sim50\,\mu$s this allows $\sim10^3$ gates within a coherence time,
the rough budget that sets near-term circuit depths.
\end{solution}
\end{exercise}

\begin{exercise}[The Kerr picture]
From \cref{eq:transmon-kerr}, find the energy of level $m$ relative to
the harmonic ladder and confirm the spacing
$\hbar\omega_{m,m+1}=\hbar\tilde\omega-mE_C$.
\begin{solution}
The Kerr term $-\tfrac{E_C}{2}\hat a^\dagger\hat a^\dagger\hat a\hat a$
has eigenvalue $-\tfrac{E_C}{2}m(m-1)$ on $\ket m$ (since
$\hat a^\dagger\hat a^\dagger\hat a\hat a\ket m=m(m-1)\ket m$). So
$E_m=\hbar\tilde\omega\,m-\tfrac{E_C}{2}m(m-1)$, and
$\hbar\omega_{m,m+1}=E_{m+1}-E_m
=\hbar\tilde\omega-\tfrac{E_C}{2}[(m+1)m-m(m-1)]
=\hbar\tilde\omega-\tfrac{E_C}{2}(2m)=\hbar\tilde\omega-mE_C$. For $m=0$
this is $\hbar\tilde\omega=\hbar\omega_{01}$, and the $0\to1$ vs $1\to2$
difference is $-E_C=\alpha$.
\end{solution}
\end{exercise}

\begin{exercise}[Why broad in phase means narrow in charge]
Explain, using $\comm{\op\varphi}{\op n}=i$, why the transmon's
charge-noise immunity is the conjugate face of its phase being pinned
in one cosine well.
\begin{solution}
Phase and number are conjugate, so a state localised in $\varphi$
(small $\varphi_{\rm zpf}$, pinned in one well) is correspondingly
\emph{delocalised} in $n$: it is a broad superposition of many charge
states. Charge dispersion arises from the wavefunction's ability to
resolve the offset $n_g$ via tunnelling between wells; a state spread
over many charge states (or equivalently, with negligible inter-well
tunnelling) cannot resolve a shift of $n_g$ by a fraction of a pair, so
its energy is insensitive to $n_g$. Large $E_J/E_C$ deepens the well,
narrows $\varphi_{\rm zpf}$, broadens the charge distribution, and
suppresses the dispersion --- all the same statement.
\end{solution}
\end{exercise}

\begin{exercise}[Reading the regime from the numbers]
Given a device with $E_J/h=12\,$GHz and (a) $E_C/h=4\,$GHz or
(b) $E_C/h=0.2\,$GHz, classify each as box-like or transmon-like and
justify.
\begin{solution}
(a) $E_J/E_C=3$: comparable energies, charge regime / intermediate ---
box-like, with substantial charge dispersion (and large
anharmonicity). (b) $E_J/E_C=60$: deep transmon regime, charge
dispersion $\propto e^{-\sqrt{480}}\approx e^{-22}$ utterly negligible,
anharmonicity $\alpha/h=-0.2\,$GHz. Device (b) is the usable qubit;
(a) would need sweet-spot operation and still dephase quickly.
\end{solution}
\end{exercise}

\begin{exercise}[Anharmonicity of the alternatives]
The transmon's anharmonicity is only $\sim5\%$. Explain qualitatively
why a flux qubit or fluxonium can achieve much larger relative
anharmonicity, and what they give up to get it.
\begin{solution}
Both place the qubit in a strongly \emph{non}-parabolic part of the
potential: the flux qubit near $\Phi_0/2$ sits in a double-well whose
two minima give well-separated, very unequal levels; the fluxonium's
shallow inductive parabola over a corrugated cosine yields levels with
large spacing differences. The transmon, by contrast, lives at the
bottom of a single deep well where the potential is nearly parabolic,
so its anharmonicity is small. The price for large anharmonicity is
sensitivity (the flux qubit to flux noise away from its sweet spot) or
complexity and low frequency (the fluxonium's superinductance and
control overhead) --- which is why the gentler transmon remains the
workhorse.
\end{solution}
\end{exercise}
